% =============================================================================
% VERZEICHNISSE (04_verzeichnisse.tex)
% =============================================================================

% Inhaltsverzeichnis
\iftoggle{showtoc}{
    \hypersetup{linkcolor=black}
    \tableofcontents
    \newpage
}{}

% Tabellenverzeichnis
\iftoggle{showtables}{
    \clearpage
    \phantomsection % Setzt den Link-Anker für Hyperref
    \addcontentsline{toc}{section}{Tabellenverzeichnis} % Erst ins TOC eintragen...
    \listoftables % ... dann das Verzeichnis generieren!
    \newpage
}{}

% Abbildungsverzeichnis
\iftoggle{showfigures}{
    \clearpage
    \phantomsection
    \addcontentsline{toc}{section}{Abbildungsverzeichnis}
    \listoffigures
    \newpage
}{}

% Oszillogrammverzeichnis
\iftoggle{showoszillos}{
    \clearpage
    \phantomsection
    \addcontentsline{toc}{section}{Oszillogrammverzeichnis}
    \listofoszillo
    \newpage
}{}

% Diagrammverzeichnis
\iftoggle{showdiagrams}{
    \clearpage
    \phantomsection
    \addcontentsline{toc}{section}{Diagrammverzeichnis}
    \listofdiagram
    \newpage
}{}


% --- ABKÜRZUNGS- UND FORMELZEICHENVERZEICHNIS ---
% (Das glossaries-Paket fügt diese durch die Option 'toc' in der Präambel bereits automatisch richtig ins TOC ein!)

% Alle Glossar-Einträge markieren, damit sie auch ohne direkte Verwendung im Text im Verzeichnis erscheinen.
\glsaddall

% Abkürzungsverzeichnis ausgeben
\printglossary[type=\acronymtype, title={Abkürzungsverzeichnis}, toctitle={Abkürzungsverzeichnis}]
\printnoidxglossary[type=\acronymtype, title={Abkürzungsverzeichnis}, toctitle={Abkürzungsverzeichnis}]
\newpage

% Formelzeichenverzeichnis ausgeben
\printglossary[type=main, title={Formelzeichenverzeichnis}, toctitle={Formelzeichenverzeichnis}]
\printnoidxglossary[type=symbols, title={Formelzeichenverzeichnis}, toctitle={Formelzeichenverzeichnis}]
\newpage
% Setzt die Linkfarbe für den Haupttext wieder auf Blau
\hypersetup{linkcolor=blue}